% ===================================================================
%  天津工业大学 · 集成电路EDA 实验报告模板 · 共用导言区
%  模板作者：陈睿晢  许可：MIT，转发、修改请保留本署名
%  版本：V-1.0.0
%
%  用法（各篇 <报告文件夹>/main.tex，路径一律从项目根目录写）：
%    \documentclass[zihao=-4,a4paper,UTF8,fontset=fandol]{ctexart}
%    % ===================================================================
%  天津工业大学 · 集成电路EDA 实验报告模板 · 共用导言区
%  模板作者：陈睿晢  许可：MIT，转发、修改请保留本署名
%  版本：V-1.0.0
%
%  用法（各篇 <报告文件夹>/main.tex，路径一律从项目根目录写）：
%    \documentclass[zihao=-4,a4paper,UTF8,fontset=fandol]{ctexart}
%    % ===================================================================
%  天津工业大学 · 集成电路EDA 实验报告模板 · 共用导言区
%  模板作者：陈睿晢  许可：MIT，转发、修改请保留本署名
%  版本：V-1.0.0
%
%  用法（各篇 <报告文件夹>/main.tex，路径一律从项目根目录写）：
%    \documentclass[zihao=-4,a4paper,UTF8,fontset=fandol]{ctexart}
%    % ===================================================================
%  天津工业大学 · 集成电路EDA 实验报告模板 · 共用导言区
%  模板作者：陈睿晢  许可：MIT，转发、修改请保留本署名
%  版本：V-1.0.0
%
%  用法（各篇 <报告文件夹>/main.tex，路径一律从项目根目录写）：
%    \documentclass[zihao=-4,a4paper,UTF8,fontset=fandol]{ctexart}
%    \input{shared/preamble.tex}
%    \renewcommand{\docdir}{<报告文件夹>}    % 本篇文件夹，\labfig 据此找图
%    \renewcommand{\labtitle}{<实验题目>}    % 以及 \classname 等报告信息
%    \begin{document}
%    \input{shared/cover.tex}
%    \input{<报告文件夹>/sections/01_objective.tex} …
%  编译器必须是 XeLaTeX（Overleaf：菜单 → Settings → Compiler），详见 README.md。
%  fontset=fandol：统一用 Fandol 字体，Overleaf 和本地编译效果一致。
% ===================================================================

% ===================== 编译器检查 =====================
% 不是 XeLaTeX 时直接停下并提示（提示用英文，避免 pdfLaTeX 下中文报错乱码）。
\RequirePackage{iftex}
\ifXeTeX\else
  \PackageError{preamble}{Compiler must be XeLaTeX. On Overleaf: Menu -> Settings -> Compiler -> XeLaTeX, then recompile. See README.md}{}
  \expandafter\stop
\fi

% ===================== 生僻字补字 =====================
% Fandol（Overleaf 默认中文字体）缺部分生僻字（如「晢」U+6662），
% 缺字时自动改用 Noto Serif CJK SC 的同一字形；其余字仍用 Fandol。
\xeCJKsetup{AutoFallBack=true}
\IfFontExistsTF{Noto Serif CJK SC}{%
  \setCJKfallbackfamilyfont{\CJKrmdefault}[BoldFont=* Bold]{Noto Serif CJK SC}%
  \setCJKfallbackfamilyfont{zhsong}[BoldFont=* Bold]{Noto Serif CJK SC}%
  \setCJKfallbackfamilyfont{zhhei}[UprightFont=* Bold]{Noto Serif CJK SC}%
  \setCJKfallbackfamilyfont{\CJKsfdefault}[UprightFont=* Bold]{Noto Serif CJK SC}%
}{}

% ===================== 宏包 =====================
\usepackage[a4paper,left=2.8cm,right=2.6cm,top=2.5cm,bottom=2.5cm]{geometry}
\usepackage{graphicx}
\usepackage{float}
\usepackage{caption}
\usepackage{subcaption}
\usepackage{booktabs}
\usepackage{enumitem}
\usepackage{xcolor}
\usepackage{listings}
\usepackage{fancyhdr}
\usepackage{setspace}
\usepackage{array}
\usepackage[hidelinks]{hyperref}

% ===================== 报告信息（各篇 main.tex 用 \renewcommand 覆盖） =====================
\newcommand{\docdir}{shared/skeleton}
\newcommand{\semester}{2026-2027学年第一学期}
\newcommand{\course}{集成电路EDA}
\newcommand{\labtitle}{}
\newcommand{\classname}{}
\newcommand{\studentid}{}
\newcommand{\studentname}{}
\newcommand{\score}{}
\newcommand{\reportdate}{2026年\qquad 月\qquad 日}

% ===================== 格式设置 =====================
\setstretch{1.4}
\setlength{\parindent}{2em}

% 一级标题：一、二、三、四、
\ctexset{
  section = {
    name = {,、},
    number = \chinese{section},
    format = \zihao{4}\heiti\raggedright,
    aftername = {},
    beforeskip = 1.2ex plus .2ex,
    afterskip = 0.8ex plus .2ex,
  },
  subsection = {
    name = {第,部分：},
    number = \chinese{subsection},
    format = \zihao{-4}\heiti\raggedright,
    aftername = {},
    beforeskip = 1ex,
    afterskip = 0.6ex,
  },
  subsubsection = {
    format = \zihao{-4}\bfseries\raggedright,
    beforeskip = 0.8ex,
    afterskip = 0.4ex,
  },
}

\captionsetup{font=small,labelsep=quad}
\captionsetup[figure]{name=图}
\captionsetup[table]{name=表}
\setlist{nosep,leftmargin=2em}

% 终端命令代码块样式
\lstset{
  basicstyle=\ttfamily\small,
  backgroundcolor=\color{gray!8},
  frame=single,
  rulecolor=\color{gray!40},
  breaklines=true,
  columns=fullflexible,
  keepspaces=true,
  showstringspaces=false,
  xleftmargin=1em,
  xrightmargin=1em,
}
\lstdefinestyle{shell}{language=bash}

% 页眉页脚
\pagestyle{fancy}
\fancyhf{}
\fancyhead[L]{\small \course\ 实验报告}
\fancyhead[R]{\small \labtitle}
\fancyfoot[C]{\small \thepage}
\renewcommand{\headrulewidth}{0.4pt}

% 插图命令：\docdir/figs/ 下图片存在则插入，否则显示占位框
% 用法：\labfig[宽度]{文件名}{图题}{标签}
\newcommand{\labfig}[4][0.85\textwidth]{%
  \begin{figure}[H]
    \centering
    \IfFileExists{\docdir/figs/#2}{%
      \includegraphics[width=#1]{\docdir/figs/#2}%
    }{%
      \fbox{\parbox[c][0.3\textheight][c]{#1}{\centering\color{gray}
        待插入截图：\texttt{figs/\detokenize{#2}}}}%
    }
    \caption{#3}
    \label{#4}
  \end{figure}
}

% 封面下划线填写栏（填写内容在下划线上居中）
\newcommand{\coverline}[2]{%
  \makebox[4em]{#1}：\underline{\makebox[16em][c]{#2}}%
}

%    \renewcommand{\docdir}{<报告文件夹>}    % 本篇文件夹，\labfig 据此找图
%    \renewcommand{\labtitle}{<实验题目>}    % 以及 \classname 等报告信息
%    \begin{document}
%    % ===================== 封面（共用；字段见 shared/preamble.tex，各篇 main.tex 覆盖） =====================
% 版式参照《实验报告封面格式.doc》：校名小一、学期四号、课程名/“实验报告”小二、信息栏小三
\begin{titlepage}
  \thispagestyle{empty}
  \centering
  \vspace*{2.5cm}

  {\zihao{1}\songti 天津工业大学\par}
  \vspace{0.6cm}
  {\zihao{4}\songti \semester\par}
  \vspace{1.8cm}

  {\zihao{-2}\heiti \course\par}
  \vspace{0.8cm}
  {\zihao{-2}\heiti 实验报告\par}
  \vspace{4cm}

  {\zihao{-3}\songti
  \renewcommand{\arraystretch}{1.9}
  \begin{tabular}{l}
    \coverline{题\hfill 目}{\labtitle} \\
    \coverline{班\hfill 级}{\classname} \\
    \coverline{学\hfill 号}{\studentid} \\
    \coverline{姓\hfill 名}{\studentname} \\
    \coverline{成\hfill 绩}{\score} \\
  \end{tabular}\par}

  \vfill
  {\zihao{-3}\songti \hfill \reportdate \hspace{3em}\par}
  \vspace{1.5cm}
\end{titlepage}

%    \input{<报告文件夹>/sections/01_objective.tex} …
%  编译器必须是 XeLaTeX（Overleaf：菜单 → Settings → Compiler），详见 README.md。
%  fontset=fandol：统一用 Fandol 字体，Overleaf 和本地编译效果一致。
% ===================================================================

% ===================== 编译器检查 =====================
% 不是 XeLaTeX 时直接停下并提示（提示用英文，避免 pdfLaTeX 下中文报错乱码）。
\RequirePackage{iftex}
\ifXeTeX\else
  \PackageError{preamble}{Compiler must be XeLaTeX. On Overleaf: Menu -> Settings -> Compiler -> XeLaTeX, then recompile. See README.md}{}
  \expandafter\stop
\fi

% ===================== 生僻字补字 =====================
% Fandol（Overleaf 默认中文字体）缺部分生僻字（如「晢」U+6662），
% 缺字时自动改用 Noto Serif CJK SC 的同一字形；其余字仍用 Fandol。
\xeCJKsetup{AutoFallBack=true}
\IfFontExistsTF{Noto Serif CJK SC}{%
  \setCJKfallbackfamilyfont{\CJKrmdefault}[BoldFont=* Bold]{Noto Serif CJK SC}%
  \setCJKfallbackfamilyfont{zhsong}[BoldFont=* Bold]{Noto Serif CJK SC}%
  \setCJKfallbackfamilyfont{zhhei}[UprightFont=* Bold]{Noto Serif CJK SC}%
  \setCJKfallbackfamilyfont{\CJKsfdefault}[UprightFont=* Bold]{Noto Serif CJK SC}%
}{}

% ===================== 宏包 =====================
\usepackage[a4paper,left=2.8cm,right=2.6cm,top=2.5cm,bottom=2.5cm]{geometry}
\usepackage{graphicx}
\usepackage{float}
\usepackage{caption}
\usepackage{subcaption}
\usepackage{booktabs}
\usepackage{enumitem}
\usepackage{xcolor}
\usepackage{listings}
\usepackage{fancyhdr}
\usepackage{setspace}
\usepackage{array}
\usepackage[hidelinks]{hyperref}

% ===================== 报告信息（各篇 main.tex 用 \renewcommand 覆盖） =====================
\newcommand{\docdir}{shared/skeleton}
\newcommand{\semester}{2026-2027学年第一学期}
\newcommand{\course}{集成电路EDA}
\newcommand{\labtitle}{}
\newcommand{\classname}{}
\newcommand{\studentid}{}
\newcommand{\studentname}{}
\newcommand{\score}{}
\newcommand{\reportdate}{2026年\qquad 月\qquad 日}

% ===================== 格式设置 =====================
\setstretch{1.4}
\setlength{\parindent}{2em}

% 一级标题：一、二、三、四、
\ctexset{
  section = {
    name = {,、},
    number = \chinese{section},
    format = \zihao{4}\heiti\raggedright,
    aftername = {},
    beforeskip = 1.2ex plus .2ex,
    afterskip = 0.8ex plus .2ex,
  },
  subsection = {
    name = {第,部分：},
    number = \chinese{subsection},
    format = \zihao{-4}\heiti\raggedright,
    aftername = {},
    beforeskip = 1ex,
    afterskip = 0.6ex,
  },
  subsubsection = {
    format = \zihao{-4}\bfseries\raggedright,
    beforeskip = 0.8ex,
    afterskip = 0.4ex,
  },
}

\captionsetup{font=small,labelsep=quad}
\captionsetup[figure]{name=图}
\captionsetup[table]{name=表}
\setlist{nosep,leftmargin=2em}

% 终端命令代码块样式
\lstset{
  basicstyle=\ttfamily\small,
  backgroundcolor=\color{gray!8},
  frame=single,
  rulecolor=\color{gray!40},
  breaklines=true,
  columns=fullflexible,
  keepspaces=true,
  showstringspaces=false,
  xleftmargin=1em,
  xrightmargin=1em,
}
\lstdefinestyle{shell}{language=bash}

% 页眉页脚
\pagestyle{fancy}
\fancyhf{}
\fancyhead[L]{\small \course\ 实验报告}
\fancyhead[R]{\small \labtitle}
\fancyfoot[C]{\small \thepage}
\renewcommand{\headrulewidth}{0.4pt}

% 插图命令：\docdir/figs/ 下图片存在则插入，否则显示占位框
% 用法：\labfig[宽度]{文件名}{图题}{标签}
\newcommand{\labfig}[4][0.85\textwidth]{%
  \begin{figure}[H]
    \centering
    \IfFileExists{\docdir/figs/#2}{%
      \includegraphics[width=#1]{\docdir/figs/#2}%
    }{%
      \fbox{\parbox[c][0.3\textheight][c]{#1}{\centering\color{gray}
        待插入截图：\texttt{figs/\detokenize{#2}}}}%
    }
    \caption{#3}
    \label{#4}
  \end{figure}
}

% 封面下划线填写栏（填写内容在下划线上居中）
\newcommand{\coverline}[2]{%
  \makebox[4em]{#1}：\underline{\makebox[16em][c]{#2}}%
}

%    \renewcommand{\docdir}{<报告文件夹>}    % 本篇文件夹，\labfig 据此找图
%    \renewcommand{\labtitle}{<实验题目>}    % 以及 \classname 等报告信息
%    \begin{document}
%    % ===================== 封面（共用；字段见 shared/preamble.tex，各篇 main.tex 覆盖） =====================
% 版式参照《实验报告封面格式.doc》：校名小一、学期四号、课程名/“实验报告”小二、信息栏小三
\begin{titlepage}
  \thispagestyle{empty}
  \centering
  \vspace*{2.5cm}

  {\zihao{1}\songti 天津工业大学\par}
  \vspace{0.6cm}
  {\zihao{4}\songti \semester\par}
  \vspace{1.8cm}

  {\zihao{-2}\heiti \course\par}
  \vspace{0.8cm}
  {\zihao{-2}\heiti 实验报告\par}
  \vspace{4cm}

  {\zihao{-3}\songti
  \renewcommand{\arraystretch}{1.9}
  \begin{tabular}{l}
    \coverline{题\hfill 目}{\labtitle} \\
    \coverline{班\hfill 级}{\classname} \\
    \coverline{学\hfill 号}{\studentid} \\
    \coverline{姓\hfill 名}{\studentname} \\
    \coverline{成\hfill 绩}{\score} \\
  \end{tabular}\par}

  \vfill
  {\zihao{-3}\songti \hfill \reportdate \hspace{3em}\par}
  \vspace{1.5cm}
\end{titlepage}

%    \input{<报告文件夹>/sections/01_objective.tex} …
%  编译器必须是 XeLaTeX（Overleaf：菜单 → Settings → Compiler），详见 README.md。
%  fontset=fandol：统一用 Fandol 字体，Overleaf 和本地编译效果一致。
% ===================================================================

% ===================== 编译器检查 =====================
% 不是 XeLaTeX 时直接停下并提示（提示用英文，避免 pdfLaTeX 下中文报错乱码）。
\RequirePackage{iftex}
\ifXeTeX\else
  \PackageError{preamble}{Compiler must be XeLaTeX. On Overleaf: Menu -> Settings -> Compiler -> XeLaTeX, then recompile. See README.md}{}
  \expandafter\stop
\fi

% ===================== 生僻字补字 =====================
% Fandol（Overleaf 默认中文字体）缺部分生僻字（如「晢」U+6662），
% 缺字时自动改用 Noto Serif CJK SC 的同一字形；其余字仍用 Fandol。
\xeCJKsetup{AutoFallBack=true}
\IfFontExistsTF{Noto Serif CJK SC}{%
  \setCJKfallbackfamilyfont{\CJKrmdefault}[BoldFont=* Bold]{Noto Serif CJK SC}%
  \setCJKfallbackfamilyfont{zhsong}[BoldFont=* Bold]{Noto Serif CJK SC}%
  \setCJKfallbackfamilyfont{zhhei}[UprightFont=* Bold]{Noto Serif CJK SC}%
  \setCJKfallbackfamilyfont{\CJKsfdefault}[UprightFont=* Bold]{Noto Serif CJK SC}%
}{}

% ===================== 宏包 =====================
\usepackage[a4paper,left=2.8cm,right=2.6cm,top=2.5cm,bottom=2.5cm]{geometry}
\usepackage{graphicx}
\usepackage{float}
\usepackage{caption}
\usepackage{subcaption}
\usepackage{booktabs}
\usepackage{enumitem}
\usepackage{xcolor}
\usepackage{listings}
\usepackage{fancyhdr}
\usepackage{setspace}
\usepackage{array}
\usepackage[hidelinks]{hyperref}

% ===================== 报告信息（各篇 main.tex 用 \renewcommand 覆盖） =====================
\newcommand{\docdir}{shared/skeleton}
\newcommand{\semester}{2026-2027学年第一学期}
\newcommand{\course}{集成电路EDA}
\newcommand{\labtitle}{}
\newcommand{\classname}{}
\newcommand{\studentid}{}
\newcommand{\studentname}{}
\newcommand{\score}{}
\newcommand{\reportdate}{2026年\qquad 月\qquad 日}

% ===================== 格式设置 =====================
\setstretch{1.4}
\setlength{\parindent}{2em}

% 一级标题：一、二、三、四、
\ctexset{
  section = {
    name = {,、},
    number = \chinese{section},
    format = \zihao{4}\heiti\raggedright,
    aftername = {},
    beforeskip = 1.2ex plus .2ex,
    afterskip = 0.8ex plus .2ex,
  },
  subsection = {
    name = {第,部分：},
    number = \chinese{subsection},
    format = \zihao{-4}\heiti\raggedright,
    aftername = {},
    beforeskip = 1ex,
    afterskip = 0.6ex,
  },
  subsubsection = {
    format = \zihao{-4}\bfseries\raggedright,
    beforeskip = 0.8ex,
    afterskip = 0.4ex,
  },
}

\captionsetup{font=small,labelsep=quad}
\captionsetup[figure]{name=图}
\captionsetup[table]{name=表}
\setlist{nosep,leftmargin=2em}

% 终端命令代码块样式
\lstset{
  basicstyle=\ttfamily\small,
  backgroundcolor=\color{gray!8},
  frame=single,
  rulecolor=\color{gray!40},
  breaklines=true,
  columns=fullflexible,
  keepspaces=true,
  showstringspaces=false,
  xleftmargin=1em,
  xrightmargin=1em,
}
\lstdefinestyle{shell}{language=bash}

% 页眉页脚
\pagestyle{fancy}
\fancyhf{}
\fancyhead[L]{\small \course\ 实验报告}
\fancyhead[R]{\small \labtitle}
\fancyfoot[C]{\small \thepage}
\renewcommand{\headrulewidth}{0.4pt}

% 插图命令：\docdir/figs/ 下图片存在则插入，否则显示占位框
% 用法：\labfig[宽度]{文件名}{图题}{标签}
\newcommand{\labfig}[4][0.85\textwidth]{%
  \begin{figure}[H]
    \centering
    \IfFileExists{\docdir/figs/#2}{%
      \includegraphics[width=#1]{\docdir/figs/#2}%
    }{%
      \fbox{\parbox[c][0.3\textheight][c]{#1}{\centering\color{gray}
        待插入截图：\texttt{figs/\detokenize{#2}}}}%
    }
    \caption{#3}
    \label{#4}
  \end{figure}
}

% 封面下划线填写栏（填写内容在下划线上居中）
\newcommand{\coverline}[2]{%
  \makebox[4em]{#1}：\underline{\makebox[16em][c]{#2}}%
}

%    \renewcommand{\docdir}{<报告文件夹>}    % 本篇文件夹，\labfig 据此找图
%    \renewcommand{\labtitle}{<实验题目>}    % 以及 \classname 等报告信息
%    \begin{document}
%    % ===================== 封面（共用；字段见 shared/preamble.tex，各篇 main.tex 覆盖） =====================
% 版式参照《实验报告封面格式.doc》：校名小一、学期四号、课程名/“实验报告”小二、信息栏小三
\begin{titlepage}
  \thispagestyle{empty}
  \centering
  \vspace*{2.5cm}

  {\zihao{1}\songti 天津工业大学\par}
  \vspace{0.6cm}
  {\zihao{4}\songti \semester\par}
  \vspace{1.8cm}

  {\zihao{-2}\heiti \course\par}
  \vspace{0.8cm}
  {\zihao{-2}\heiti 实验报告\par}
  \vspace{4cm}

  {\zihao{-3}\songti
  \renewcommand{\arraystretch}{1.9}
  \begin{tabular}{l}
    \coverline{题\hfill 目}{\labtitle} \\
    \coverline{班\hfill 级}{\classname} \\
    \coverline{学\hfill 号}{\studentid} \\
    \coverline{姓\hfill 名}{\studentname} \\
    \coverline{成\hfill 绩}{\score} \\
  \end{tabular}\par}

  \vfill
  {\zihao{-3}\songti \hfill \reportdate \hspace{3em}\par}
  \vspace{1.5cm}
\end{titlepage}

%    \input{<报告文件夹>/sections/01_objective.tex} …
%  编译器必须是 XeLaTeX（Overleaf：菜单 → Settings → Compiler），详见 README.md。
%  fontset=fandol：统一用 Fandol 字体，Overleaf 和本地编译效果一致。
% ===================================================================

% ===================== 编译器检查 =====================
% 不是 XeLaTeX 时直接停下并提示（提示用英文，避免 pdfLaTeX 下中文报错乱码）。
\RequirePackage{iftex}
\ifXeTeX\else
  \PackageError{preamble}{Compiler must be XeLaTeX. On Overleaf: Menu -> Settings -> Compiler -> XeLaTeX, then recompile. See README.md}{}
  \expandafter\stop
\fi

% ===================== 生僻字补字 =====================
% Fandol（Overleaf 默认中文字体）缺部分生僻字（如「晢」U+6662），
% 缺字时自动改用 Noto Serif CJK SC 的同一字形；其余字仍用 Fandol。
\xeCJKsetup{AutoFallBack=true}
\IfFontExistsTF{Noto Serif CJK SC}{%
  \setCJKfallbackfamilyfont{\CJKrmdefault}[BoldFont=* Bold]{Noto Serif CJK SC}%
  \setCJKfallbackfamilyfont{zhsong}[BoldFont=* Bold]{Noto Serif CJK SC}%
  \setCJKfallbackfamilyfont{zhhei}[UprightFont=* Bold]{Noto Serif CJK SC}%
  \setCJKfallbackfamilyfont{\CJKsfdefault}[UprightFont=* Bold]{Noto Serif CJK SC}%
}{}

% ===================== 宏包 =====================
\usepackage[a4paper,left=2.8cm,right=2.6cm,top=2.5cm,bottom=2.5cm]{geometry}
\usepackage{graphicx}
\usepackage{float}
\usepackage{caption}
\usepackage{subcaption}
\usepackage{booktabs}
\usepackage{enumitem}
\usepackage{xcolor}
\usepackage{listings}
\usepackage{fancyhdr}
\usepackage{setspace}
\usepackage{array}
\usepackage[hidelinks]{hyperref}

% ===================== 报告信息（各篇 main.tex 用 \renewcommand 覆盖） =====================
\newcommand{\docdir}{shared/skeleton}
\newcommand{\semester}{2026-2027学年第一学期}
\newcommand{\course}{集成电路EDA}
\newcommand{\labtitle}{}
\newcommand{\classname}{}
\newcommand{\studentid}{}
\newcommand{\studentname}{}
\newcommand{\score}{}
\newcommand{\reportdate}{2026年\qquad 月\qquad 日}

% ===================== 格式设置 =====================
\setstretch{1.4}
\setlength{\parindent}{2em}

% 一级标题：一、二、三、四、
\ctexset{
  section = {
    name = {,、},
    number = \chinese{section},
    format = \zihao{4}\heiti\raggedright,
    aftername = {},
    beforeskip = 1.2ex plus .2ex,
    afterskip = 0.8ex plus .2ex,
  },
  subsection = {
    name = {第,部分：},
    number = \chinese{subsection},
    format = \zihao{-4}\heiti\raggedright,
    aftername = {},
    beforeskip = 1ex,
    afterskip = 0.6ex,
  },
  subsubsection = {
    format = \zihao{-4}\bfseries\raggedright,
    beforeskip = 0.8ex,
    afterskip = 0.4ex,
  },
}

\captionsetup{font=small,labelsep=quad}
\captionsetup[figure]{name=图}
\captionsetup[table]{name=表}
\setlist{nosep,leftmargin=2em}

% 终端命令代码块样式
\lstset{
  basicstyle=\ttfamily\small,
  backgroundcolor=\color{gray!8},
  frame=single,
  rulecolor=\color{gray!40},
  breaklines=true,
  columns=fullflexible,
  keepspaces=true,
  showstringspaces=false,
  xleftmargin=1em,
  xrightmargin=1em,
}
\lstdefinestyle{shell}{language=bash}

% 页眉页脚
\pagestyle{fancy}
\fancyhf{}
\fancyhead[L]{\small \course\ 实验报告}
\fancyhead[R]{\small \labtitle}
\fancyfoot[C]{\small \thepage}
\renewcommand{\headrulewidth}{0.4pt}

% 插图命令：\docdir/figs/ 下图片存在则插入，否则显示占位框
% 用法：\labfig[宽度]{文件名}{图题}{标签}
\newcommand{\labfig}[4][0.85\textwidth]{%
  \begin{figure}[H]
    \centering
    \IfFileExists{\docdir/figs/#2}{%
      \includegraphics[width=#1]{\docdir/figs/#2}%
    }{%
      \fbox{\parbox[c][0.3\textheight][c]{#1}{\centering\color{gray}
        待插入截图：\texttt{figs/\detokenize{#2}}}}%
    }
    \caption{#3}
    \label{#4}
  \end{figure}
}

% 封面下划线填写栏（填写内容在下划线上居中）
\newcommand{\coverline}[2]{%
  \makebox[4em]{#1}：\underline{\makebox[16em][c]{#2}}%
}
